\documentclass[12pt, ukrainian]{article}
\usepackage[a4paper]{geometry}
\setlength\parindent{0mm}
\geometry{top=1.1cm,bottom=1.1cm,left=1.1cm,right=1.1cm}
\pagestyle{empty}
\usepackage[T2A]{fontenc}
\usepackage[utf8]{inputenc}
\usepackage[ukrainian]{babel}
\usepackage{graphicx}
\usepackage{caption}
\usepackage{verbatim, listings, clrscode3e, fancyvrb}
\DeclareCaptionFormat{nocaption}{}
\captionsetup{format=nocaption,aboveskip=0pt,belowskip=0pt}
\usepackage[Export]{adjustbox}
\adjustboxset{max size={0.9\linewidth}{0.9\paperheight}}
\def\code#1{\Verb!#1!}
\begin{document}
{{21.05.2025 Контрольна робота, приклад } \par}
{
%\begin {large}
У завданнях 1-7 вкажіть, що буде виведено за виконання виклику test(); . 
    Вважати, що всі необхідні включення, декларації тощо наявні. 
    У випадку будь-якої помилки (синтаксичної, часу виконання, невизначеної поведінки тощо) зазначте, наявність помилки та поясніть у чому саме помилка.
\begin{minipage}[t]{11cm}
               
{\begin{center}Завдання {1}\end{center}}
\vspace{-6mm}

\begin{verbatim}
int z=-3;
struct A{
  A(){z+=3;}
  A(A &aa){z++;}
  ~A(){z++;}
};
void test(){
  {A a, b(a), c(b);} cout<<z;}
\end{verbatim}            


               
{\begin{center}Завдання {2}\end{center}}
\vspace{-6mm}

\begin{verbatim}
struct A{ A(){cout<<"A";}};
struct B{ B(){cout<<"B";}};
class C{
  B b[1]; A a[3];
public:
  C(){cout<<"C";} };
void test(){C c;}
\end{verbatim}            


               
{\begin{center}Завдання {3}\end{center}}
\vspace{-6mm}

\begin{verbatim}
struct A{
  A(){cout<<7;}
  A(const A&){cout<<6;}
  A(A&&){cout<<8;}
  A& operator=(const A&){cout<<3;return *this;}
  A& operator=(A&&){cout<<2;return *this;}};
void test(){A a, b(a), c(move(a)), d=a;}
\end{verbatim}            


               
{\begin{center}Завдання {5}\end{center}}
\vspace{-6mm}

\begin{verbatim}
struct A{
  virtual int work(){return 1;}
  int relax(){return 6;}};
struct B:public A{ int relax(){return 5;}};
struct C:public B{virtual int relax(){return 2;}};
void test(){ A *a1=new B, *a2=new C;
  cout<<a1->work()<<a2->relax()<<a1->relax();}
\end{verbatim}            


\end{minipage}
\vline \hspace{8mm}
\begin{minipage}[t]{7cm}
               
{\begin{center}Завдання {4}\end{center}}
\vspace{-6mm}

\begin{verbatim}
struct A{
  int x;
  virtual void d(){x+=3;}
  void g(){A::d();}};
struct B:public A{
  int y;
  void d(){y-=3;}};
void test(){
  B b; b.x=b.y=4; b.g();
  cout<<b.x<<b.y;}
\end{verbatim}            


               
{\begin{center}Завдання {6}\end{center}}
\vspace{-6mm}

\begin{verbatim}
class A{};
void test(){
  A *a=new A[2]; 
  ++a; delete a; cout<<"=";} 
\end{verbatim}            


               
{\begin{center}Завдання {7}\end{center}}
\vspace{-6mm}

\begin{verbatim}
struct A{
  A(){throw 1;cout<<"A";}
  ~A(){cout<<"~A";}};
struct B{
  B(){throw 2;cout<<"B";}
  ~B(){cout<<"~B";}};
class C{B b; A a;
public:
  C(){cout<<"C";}
  ~C(){cout<<"~C";}};
void test(){
  try{C c;} 
  catch(...){} 
  cout<<"=";}
\end{verbatim}            


\end{minipage}
               
{\begin{center}Завдання {8}\end{center}}
\vspace{-6mm}
Задано тип вузла списку
struct Node \{int n; Node *prev, *next;\};
У списку зберігаються послідовні цілі числа від -24 до 40. Вказівник p встановлено на вузол, в якому зберігається число 3.
Чому дорівнює значення p->next->next->prev->prev->next->n ?\par

               
{\begin{center}Завдання {9}\end{center}}
\vspace{-6mm}
В умовах попередньої задачі було виконано p->next=p->prev->prev;

Чому дорівнює значення p->next->next->prev->next->n ?\par

\newpage
               
{\begin{center}Завдання {10}\end{center}}
\vspace{-6mm}
На початку структура даних стек порожня. Далі над нею було виконано операції:

додати елемент~1, додати елемент~2, вилучити елемент, додати елемент~4, вилучити елемент, додати елемент~6, вилучити елемент, додати елемент~7, додати елемент~8, вилучити елемент

Який елемент було вилучено останнім?\par

               
{\begin{center}Завдання {11}\end{center}}
\vspace{-6mm}

 Написати функцію ...
Якість коду також оцінюється.\par

%end
\end{document}